\documentclass{amsart}
% A tiny paper written for Madalean's tests (CC0): labels, references, a proof citing two lemmas,
% a « Proof of Theorem » section, a remark that is not a result, and a proposition that cites nothing.
\newtheorem{theorem}{Theorem}[section]
\newtheorem{lemma}[theorem]{Lemma}
\newtheorem{proposition}[theorem]{Proposition}
\theoremstyle{definition}
\newtheorem{definition}[theorem]{Definition}
\newtheorem*{remark}{Remark}
\title{Small facts about natural numbers}

\begin{document}
\maketitle

\section{Preliminaries}

\begin{definition}\label{def:double}
The double of a natural number $n$ is $n + n$.
\end{definition}

\begin{lemma}\label{lem:assoc}
For all natural numbers $a, b, c$ we have $a + b + c = a + (b + c)$.
\end{lemma}

\begin{lemma}[Commutativity]\label{lem:comm}
For all natural numbers $a$ and $b$ we have $a + b = b + a$.
\end{lemma}

\begin{lemma}\label{lem:square}
For all natural numbers $a$ and $b$ we have $(a + b)(a + b) = a a + b (2 a + b)$.
\end{lemma}
\begin{proof}
Expand the product and use Lemma~\ref{lem:comm}.
\end{proof}

\section{Main results}

\begin{theorem}\label{thm:main}
For all natural numbers $a, b, c$ we have $(a + b)(a + b) + c = a a + (b (2 a + b) + c)$.
\end{theorem}
\begin{proof}
By Lemma~\ref{lem:square} and Lemma~\ref{lem:assoc}, as in \cite{euclid}.
\end{proof}

\begin{remark}
Theorem~\ref{thm:main} also holds in every commutative monoid.
\end{remark}

\begin{theorem}\label{thm:late}
For every natural number $n$, the double of $n$ is $2 n$.
\end{theorem}

\begin{proposition}\label{prop:alone}
Every natural number is less than its successor.
\end{proposition}

\section{Proof of Theorem~\ref{thm:late}}
By Definition~\ref{def:double} the double of $n$ is $n + n$, and we only use that $2 n = n + n$.

\begin{thebibliography}{9}
\bibitem{euclid} Euclid, \emph{Elements}.
\end{thebibliography}
\end{document}
